\documentclass[11pt,a4paper]{article}
\usepackage[utf8]{inputenc}
\usepackage[spanish]{babel}
\usepackage{amsmath}
\usepackage{amssymb}
\usepackage{geometry}
\geometry{margin=1in}

\title{Demostración: Ortogonalidad, Normalización y Relaciones de Recurrencia de las Funciones de Bessel en el marco de Sturm-Liouville}
\author{}
\date{}

\begin{document}

\maketitle

\section{Introducción y Planteamiento}
El objetivo es demostrar la relación de ortogonalidad y normalización para las funciones de Bessel de primera especie $J_n$, dada por:
\begin{equation}
	\int_0^a J_n(k_{nj}r/a)J_n(k_{ni}r/a)rdr = \frac{a^2}{2}\delta_{ij} [J'_n(k_{ni})]^2 = \frac{a^2}{2}\delta_{ij}[J_{n+1}(k_{ni})]^2
\end{equation}
donde $k_{ni}$ y $k_{nj}$ son los ceros de la función de Bessel $J_n(x)$, es decir, $J_n(k_{ni}) = 0$ y $J_n(k_{nj}) = 0$.

Para simplificar la notación durante el desarrollo, definiremos las constantes:
\begin{equation}
	\alpha = \frac{k_{ni}}{a}, \quad \beta = \frac{k_{nj}}{a}
\end{equation}
De este modo, las funciones bajo estudio son $y_i(r) = J_n(\alpha r)$ y $y_j(r) = J_n(\beta r)$, con las condiciones de frontera $y_i(a) = y_j(a) = 0$.

\section{El Marco de Sturm-Liouville y la Ortogonalidad ($i \neq j$)}
La ecuación diferencial de Bessel para $y(r) = J_n(\alpha r)$ se expresa ordinariamente como:
\begin{equation}
	r^2 y'' + r y' + (\alpha^2 r^2 - n^2) y = 0
\end{equation}
Dividiendo por $r$, podemos escribirla en la forma auto-adjunta o forma estándar de Sturm-Liouville, $[p(r)y']' + [q(r) + \lambda w(r)]y = 0$:
\begin{equation}
	\frac{d}{dr}\left( r \frac{dy}{dr} \right) + \left( -\frac{n^2}{r} + \alpha^2 r \right) y = 0
\end{equation}
A partir de esta forma, identificamos los elementos del sistema de Sturm-Liouville:
\begin{itemize}
	\item Función de ponderación (peso): $w(r) = r$
	\item Coeficiente principal: $p(r) = r$
	\item Función potencial: $q(r) = -\frac{n^2}{r}$
	\item Parámetro del eigenvalor: $\lambda = \alpha^2$
\end{itemize}

La teoría general de Sturm-Liouville establece que para un operador auto-adjunto con condiciones de frontera homogéneas en un intervalo $[0, a]$, las eigenfunciones correspondientes a eigenvalores distintos ($\alpha \neq \beta$, lo cual implica $i \neq j$) son ortogonales respecto a la función de peso $w(r)$. Por lo tanto:
\begin{equation}
	\int_0^a J_n\left(\frac{k_{nj}}{a}r\right) J_n\left(\frac{k_{ni}}{a}r\right) r \, dr = 0 \quad \text{para } i \neq j
\end{equation}
Este resultado valida la presencia de la delta de Kronecker $\delta_{ij}$ en la expresión general.

\section{Condición de Normalización ($i = j$)}
Cuando $i = j$, se tiene $\alpha = \beta$. Para evaluar la integral del cuadrado de la función, partimos de la ecuación diferencial en forma de Sturm-Liouville para $y(r) = J_n(\alpha r)$:
\begin{equation}
	\frac{d}{dr}(r y') + \left(\alpha^2 r - \frac{n^2}{r}\right) y = 0
\end{equation}
Multiplicamos toda la ecuación por el factor $2r y'$:
\begin{equation}
	2r y' \frac{d}{dr}(r y') + 2\alpha^2 r^2 y y' - 2n^2 y y' = 0
\end{equation}
Reconocemos que cada uno de estos términos representa una derivada total respecto a $r$:
\begin{align}
	2r y' \frac{d}{dr}(r y') & = \frac{d}{dr} \left[ r^2 (y')^2 \right] \\
	2\alpha^2 r^2 y y'       & = \alpha^2 r^2 \frac{d}{dr} (y^2)        \\
	2n^2 y y'                & = n^2 \frac{d}{dr} (y^2)
\end{align}
Sustituyendo estas identidades, obtenemos:
\begin{equation}
	\frac{d}{dr} \left[ r^2 (y')^2 \right] + (\alpha^2 r^2 - n^2) \frac{d}{dr} (y^2) = 0
\end{equation}
Procedemos a integrar ambos miembros en el intervalo $[0, a]$:
\begin{equation}
	\int_0^a \frac{d}{dr} \left[ r^2 (y')^2 \right] dr + \int_0^a (\alpha^2 r^2 - n^2) \frac{d}{dr} (y^2) dr = 0
\end{equation}
El primer término se integra de forma directa por el Teorema Fundamental del Cálculo. Para el segundo término, aplicamos integración por partes definiendo $u = \alpha^2 r^2 - n^2$ y $dv = \frac{d}{dr}(y^2) dr$, lo que implica $du = 2\alpha^2 r \, dr$ y $v = y^2$:
\begin{equation}
	\left[ r^2 (y')^2 \right]_0^a + \left[ (\alpha^2 r^2 - n^2) y^2 \right]_0^a - \int_0^a y^2 (2\alpha^2 r) \, dr = 0
\end{equation}
Evaluamos los límites de integración:
\begin{itemize}
	\item En $r = a$: Por hipótesis, $y(a) = J_n(\alpha a) = J_n(k_{ni}) = 0$. Esto anula por completo el segundo término evaluado en el límite superior. El primer término nos proporciona $a^2 [y'(a)]^2$.
	\item En $r = 0$: El término $r^2 (y')^2$ se desvanece debido al factor $r^2$ (dado que $y'(0)$ es finito para todo $n \ge 0$). Asimismo, para el segundo término, si $n > 0$, $y(0) = 0$; si $n = 0$, el factor $(\alpha^2 r^2 - n^2)$ se reduce a $0$ en $r=0$. En cualquier caso, el límite inferior es nulo.
\end{itemize}
De este modo, la ecuación se reduce a:
\begin{equation}
	a^2 [y'(a)]^2 - 2\alpha^2 \int_0^a r [J_n(\alpha r)]^2 dr = 0
\end{equation}
Despejando la integral de interés:
\begin{equation}
	\int_0^a r [J_n(\alpha r)]^2 dr = \frac{a^2}{2\alpha^2} [y'(a)]^2
\end{equation}
Aplicando la regla de la cadena a la función original $y(r) = J_n(\alpha r)$, su derivada es $y'(r) = \alpha J'_n(\alpha r)$. Al evaluar en $r = a$:
\begin{equation}
	y'(a) = \alpha J'_n(\alpha a) = \alpha J'_n(k_{ni})
\end{equation}
Sustituyendo este resultado en la expresión de la integral:
\begin{equation}
	\int_0^a r [J_n(\alpha r)]^2 dr = \frac{a^2}{2\alpha^2} \left[ \alpha J'_n(k_{ni}) \right]^2 = \frac{a^2}{2} [J'_n(k_{ni})]^2
\end{equation}
Combinando esto con la condición de ortogonalidad obtenida en la sección anterior, queda demostrada la primera igualdad:
\begin{equation}
	\int_0^a J_n(k_{nj}r/a)J_n(k_{ni}r/a)rdr = \frac{a^2}{2}\delta_{ij} [J'_n(k_{ni})]^2
\end{equation}

\section{Deducción de la Relación de Recurrencia}
Para fundamentar rigurosamente la transición a la última igualdad, deduciremos la relación diferencial de recurrencia a partir de la definición en serie de potencias de las funciones de Bessel de primera especie de orden $n$:
\begin{equation}
	J_n(x) = \sum_{m=0}^{\infty} \frac{(-1)^m}{m!(m+n)!} \left(\frac{x}{2}\right)^{2m+n}
\end{equation}
Derivando término a término respecto a $x$ mediante la regla de la cadena, se obtiene:
\begin{equation}
	J'_n(x) = \sum_{m=0}^{\infty} \frac{(-1)^m}{m!(m+n)!} (2m+n) \left( \frac{1}{2} \right) \left(\frac{x}{2}\right)^{2m+n-1}
\end{equation}
El factor lineal $(2m+n)$ dentro del sumatorio se puede separar de forma aditiva en dos sumas distintas:
\begin{equation}
	J'_n(x) = \sum_{m=0}^{\infty} \frac{(-1)^m n}{2 \cdot m!(m+n)!} \left(\frac{x}{2}\right)^{2m+n-1} + \sum_{m=0}^{\infty} \frac{(-1)^m 2m}{2 \cdot m!(m+n)!} \left(\frac{x}{2}\right)^{2m+n-1}
\end{equation}

Analicemos la primera sumatoria de la expresión anterior. Multiplicando y dividiendo convenientemente por $x$, podemos reescribir el exponente para recuperar la forma original de $J_n(x)$:
\begin{equation}
	\frac{n}{2} \left(\frac{x}{2}\right)^{2m+n-1} = \frac{n}{x} \left(\frac{x}{2}\right) \left(\frac{x}{2}\right)^{2m+n-1} = \frac{n}{x} \left(\frac{x}{2}\right)^{2m+n}
\end{equation}
Extrayendo el factor común $\frac{n}{x}$ fuera del operador de suma, identificamos plenamente la serie de $J_n(x)$:
\begin{equation}
	\text{Suma}_1 = \frac{n}{x} \sum_{m=0}^{\infty} \frac{(-1)^m}{m!(m+n)!} \left(\frac{x}{2}\right)^{2m+n} = \frac{n}{x} J_n(x)
\end{equation}

Para la segunda sumatoria, simplificamos el factor $\frac{2m}{2} = m$. Puesto que el término correspondiente a $m=0$ se anula, modificamos el límite inferior del índice a $m=1$. Esto nos permite cancelar la variable $m$ con el primer componente del factorial del denominador, utilizando la propiedad $m! = m(m-1)!$:
\begin{equation}
	\text{Suma}_2 = \sum_{m=1}^{\infty} \frac{(-1)^m}{(m-1)!(m+n)!} \left(\frac{x}{2}\right)^{2m+n-1}
\end{equation}
Efectuamos un cambio de variable para los índices definiendo $k = m - 1$ (de donde $m = k + 1$). Bajo esta transformación:
\begin{itemize}
	\item El límite inferior $m=1$ se convierte en $k=0$.
	\item El factor alternante $(-1)^m$ pasa a ser $(-1)^{k+1} = -(-1)^k$.
	\item El factorial $(m+n)!$ se reescribe como $(k+n+1)!$.
	\item El exponente de la potencia se transforma en $2(k+1)+n-1 = 2k+(n+1)$.
\end{itemize}
Sustituyendo estos nuevos términos en la sumatoria:
\begin{equation}
	\text{Suma}_2 = -\sum_{k=0}^{\infty} \frac{(-1)^k}{k!(k+n+1)!} \left(\frac{x}{2}\right)^{2k+(n+1)}
\end{equation}
Esta última estructura matemática coincide de manera exacta con la definición en serie de la función de Bessel de orden $n+1$. Por consiguiente:
\begin{equation}
	\text{Suma}_2 = -J_{n+1}(x)
\end{equation}

Combinando los resultados obtenidos para ambas sumatorias ($\text{Suma}_1 + \text{Suma}_2$), se demuestra la identidad fundamental de recurrencia diferencial:
\begin{equation}
	J'_n(x) = \frac{n}{x} J_n(x) - J_{n+1}(x)
\end{equation}

\section{Finalización de la Demostración del Ejercicio}
Una vez establecida rigurosamente la relación de recurrencia, procedemos a evaluarla específicamente en los puntos críticos dados por las raíces, es decir, $x = k_{ni}$:
\begin{equation}
	J'_n(k_{ni}) = \frac{n}{k_{ni}} J_n(k_{ni}) - J_{n+1}(k_{ni})
\end{equation}
Haciendo uso de la condición de frontera fundamental del problema, $J_n(k_{ni}) = 0$, el primer término de la derecha se anula idénticamente, dando lugar a:
\begin{equation}
	J'_n(k_{ni}) = -J_{n+1}(k_{ni})
\end{equation}
Al elevar al cuadrado ambos miembros de la ecuación, el signo negativo se disipa por propiedades algebraicas elementales:
\begin{equation}
	[J'_n(k_{ni})]^2 = [J_{n+1}(k_{ni})]^2
\end{equation}
Sustituyendo finalmente este resultado en la integral de normalización obtenida en la Sección 3, se alcanza el resultado definitivo plasmado en el enunciado original:
\begin{equation}
	\int_0^a J_n(k_{nj}r/a)J_n(k_{ni}r/a)rdr = \frac{a^2}{2}\delta_{ij}[J_{n+1}(k_{ni})]^2
\end{equation}
Q.E.D.

\end{document}
